\documentclass{beamer}
\usepackage[utf8]{inputenc}
\usepackage[spanish]{babel}
\usepackage{graphicx}
\usepackage{booktabs}
\usepackage{hyperref}

\usetheme{Madrid}
\usecolortheme{beaver}

% ---------------------------------------------------
% PORTADA
% ---------------------------------------------------
\title[Festival de las Luciérnagas]{Sistema Integral de Reservaciones: Festival de las Luciérnagas}
\subtitle{Proyecto de Desarrollo de Software Seguro y Escalable}
\author[Equipo de Desarrollo]{Empresa: LuciDev Solutions \\ \vspace{0.3cm} Integrantes: \\ Zassh \and Equipo de Desarrollo}
\institute[LuciDev]{Universidad / Instituto}
\date{\today}

\begin{document}

\begin{frame}
  \titlepage
\end{frame}

\begin{frame}{Índice}
  \tableofcontents
\end{frame}

% ---------------------------------------------------
% REQUERIMIENTOS
% ---------------------------------------------------
\section{Requerimientos}
\begin{frame}{Requerimientos del Sistema}
  \textbf{Requerimientos Funcionales:}
  \begin{itemize}
      \item Registro y autenticación de usuarios (Clientes y Administradores).
      \item Visualización de santuarios con mapas interactivos.
      \item Sistema de reservas de cabañas y zonas de camping.
      \item Generación de Pases de Acceso en PDF con código QR.
      \item Dashboard de administrador para validar ingresos mediante escáner QR.
  \end{itemize}
  \vspace{0.3cm}
  \textbf{Requerimientos No Funcionales:}
  \begin{itemize}
      \item \textbf{Seguridad:} Contraseñas encriptadas, protección de rutas y prevención de inyecciones.
      \item \textbf{Usabilidad:} Interfaz intuitiva y estética con modo oscuro (Premium Design).
      \item \textbf{Rendimiento:} Tiempos de carga menores a 2 segundos mediante compilación optimizada (Vite/React).
  \end{itemize}
\end{frame}

% ---------------------------------------------------
% PROPUESTA DE VALOR
% ---------------------------------------------------
\section{Propuesta de Valor}
\begin{frame}{Propuesta de Valor del Proyecto}
  \begin{block}{Innovación Tecnológica y Ecológica}
    El \textbf{Festival de las Luciérnagas} ofrece una plataforma digital que moderniza la gestión turística de áreas naturales protegidas. 
  \end{block}
  \vspace{0.2cm}
  Nuestra propuesta radica en:
  \begin{itemize}
      \item \textbf{Digitalización Cero Papel:} Entradas 100\% digitales mediante códigos QR y PDFs generados on-the-fly.
      \item \textbf{Control de Aforo Seguro:} Evita la sobreexplotación de los santuarios al manejar inventarios de reservas en tiempo real.
      \item \textbf{Experiencia Premium:} Un diseño de interfaz inmersivo que captura la magia de las luciérnagas antes de llegar al bosque.
  \end{itemize}
\end{frame}

% ---------------------------------------------------
% CICLO DE VIDA
% ---------------------------------------------------
\section{Ciclo de Vida}
\begin{frame}{Ciclo de Vida del Software}
  Se adoptó un \textbf{Modelo Iterativo e Incremental}, permitiendo entregas funcionales continuas.
  \vspace{0.3cm}
  \begin{enumerate}
      \item \textbf{Análisis:} Identificación del problema de sobrecupo en santuarios.
      \item \textbf{Diseño:} Creación de prototipos UI/UX y modelado de base de datos relacional.
      \item \textbf{Implementación:} Desarrollo del backend (Node.js/Prisma) y frontend (React).
      \item \textbf{Pruebas:} Validación de roles, autenticación y escaneo de QR.
      \item \textbf{Despliegue:} Configuración en VPS (Ubuntu) utilizando PM2 y Nginx.
      \item \textbf{Mantenimiento:} Refactorizaciones basadas en la retroalimentación.
  \end{enumerate}
\end{frame}

% ---------------------------------------------------
% METODOLOGÍA
% ---------------------------------------------------
\section{Metodología}
\begin{frame}{Metodología(s) Utilizada(s)}
  Se implementó \textbf{Scrum} adaptado para equipos ágiles pequeños:
  \begin{itemize}
      \item \textbf{Sprints:} Ciclos de desarrollo de 1 a 2 semanas.
      \item \textbf{Product Backlog:} Historias de usuario claras (e.g., "Como admin, quiero escanear un QR para dar acceso").
      \item \textbf{Daily Standups:} Revisiones de progreso rápidas.
  \end{itemize}
  \vspace{0.3cm}
  Prácticas Complementarias:
  \begin{itemize}
      \item \textbf{Pair Programming:} Especialmente durante la integración de pasarelas y seguridad.
      \item \textbf{Control de Versiones:} Uso de Git/GitHub con flujos de trabajo basados en ramas (Feature Branches).
  \end{itemize}
\end{frame}

% ---------------------------------------------------
% SEGURIDAD
% ---------------------------------------------------
\section{Seguridad}
\begin{frame}{Aspectos de Seguridad Considerados}
  La plataforma implementa estándares modernos de seguridad web:
  \vspace{0.2cm}
  \begin{itemize}
      \item \textbf{Autenticación y Autorización:} Uso de JSON Web Tokens (JWT) almacenados de forma segura y validación estricta de Roles (RBAC).
      \item \textbf{Protección de Datos:} Encriptación de contraseñas usando \texttt{bcrypt} (costo factor 12).
      \item \textbf{Protección contra Vulnerabilidades Comunes:}
      \begin{itemize}
          \item Sanitización de entradas (prevención de XSS y SQL Injection) a través del ORM Prisma y validadores como Zod.
          \item CORS configurado estrictamente.
      \end{itemize}
      \item \textbf{Restricciones de Hardware:} Manejo de errores gracefully en la API del escáner de cámara bajo contextos no-HTTPS.
  \end{itemize}
\end{frame}

% ---------------------------------------------------
% REFACTORIZACIONES Y CLEAN CODE
% ---------------------------------------------------
\section{Código Limpio}
\begin{frame}{Refactorizaciones y Código Limpio}
  El proyecto mantiene altos estándares de calidad de código:
  \begin{block}{Principios Aplicados}
      \begin{itemize}
          \item \textbf{DRY (Don't Repeat Yourself):} Creación de Hooks personalizados (\texttt{useAuth}, \texttt{useParks}) y componentes reutilizables (Botones, Tarjetas).
          \item \textbf{Separation of Concerns:} Frontend separado del Backend. Backend estructurado en arquitectura por capas (Controladores, Servicios, Rutas).
          \item \textbf{Tipado Fuerte:} Uso extensivo de TypeScript para prevenir errores en tiempo de ejecución.
      \end{itemize}
  \end{block}
  \vspace{0.2cm}
  \textbf{Ejemplo de Refactorización:} 
  Se reemplazó una tabla genérica de reservas por un diseño de "Tarjetas de Pase" con descarga directa a PDF (\texttt{jsPDF}), separando la lógica de renderizado del manejo de archivos.
\end{frame}

% ---------------------------------------------------
% IMPLEMENTACIÓN
% ---------------------------------------------------
\section{Implementación}
\begin{frame}{Muestra de la Implementación}
  \begin{center}
      \textit{(En la presentación en vivo, realizar la demostración del software aquí)}
  \end{center}
  \vspace{0.2cm}
  \textbf{Puntos clave a demostrar:}
  \begin{enumerate}
      \item Inicio de sesión con cuenta de Administrador.
      \item Interfaz gráfica y renderizado del mapa.
      \item Generación de una reserva de prueba.
      \item Visualización de la reserva en el perfil del cliente.
      \item Descarga del Pase Mágico en formato PDF.
      \item Escaneo del código QR (o ingreso manual del ID) en el panel de administrador para confirmar la llegada.
  \end{enumerate}
\end{frame}

% ---------------------------------------------------
% BITÁCORA
% ---------------------------------------------------
\section{Bitácora}
\begin{frame}{Bitácora de Trabajo en Equipo}
  \begin{table}
      \centering
      \resizebox{\textwidth}{!}{%
      \begin{tabular}{@{}llp{6cm}@{}}
      \toprule
      \textbf{Fecha} & \textbf{Fase} & \textbf{Actividad Principal} \\ \midrule
      Semana 1 & Planificación & Definición de requerimientos y diseño de base de datos (Prisma). \\
      Semana 2 & Backend & Desarrollo de la API REST, autenticación JWT y roles. \\
      Semana 3 & Frontend UI & Maquetación con Tailwind CSS, integración de Framer Motion. \\
      Semana 4 & Integración & Conexión Frontend-Backend, validación de estado de reservas. \\
      Semana 5 & Refinamiento & Generación de PDF (jsPDF), escáner QR y fixes de HTTPS. \\
      Semana 6 & Despliegue & Puesta en producción en Servidor VPS (Nginx + PM2). \\ \bottomrule
      \end{tabular}%
      }
  \end{table}
\end{frame}

% ---------------------------------------------------
% BIBLIOGRAFÍA
% ---------------------------------------------------
\section{Bibliografía}
\begin{frame}{Bibliografía (Formato APA)}
  \begin{thebibliography}{9}
      \bibitem{react} 
      Meta Platforms, Inc. (2024). \textit{React: The library for web and native user interfaces}. Recuperado de https://react.dev/
      
      \bibitem{prisma} 
      Prisma Data, Inc. (2024). \textit{Prisma: Next-generation Node.js and TypeScript ORM}. Recuperado de https://www.prisma.io/
      
      \bibitem{tailwind} 
      Wathan, A. (2024). \textit{Tailwind CSS: Rapidly build modern websites without ever leaving your HTML}. Recuperado de https://tailwindcss.com/
      
      \bibitem{martin}
      Martin, R. C. (2008). \textit{Clean Code: A Handbook of Agile Software Craftsmanship}. Prentice Hall.
      
      \bibitem{sommerville}
      Sommerville, I. (2015). \textit{Ingeniería de Software} (10ma ed.). Pearson.
  \end{thebibliography}
\end{frame}

\begin{frame}
  \begin{center}
      \Huge \textbf{¡Gracias!} \\
      \vspace{0.5cm}
      \Large ¿Preguntas?
  \end{center}
\end{frame}

\end{document}
